\begin{authorclaim}{Theorem4}
For every $\epsilon>0$ there exists a palette-alternating tree code with rate $1-\epsilon$ and distance $\delta=\Omega(\epsilon\cdot\log^{-1}(1/\epsilon))$.
\end{authorclaim}
\paragraph{Comparison with the Gilbert-Varshamov bound}
Observe that the distance-rate trade-off obtained in Theorem4 is the same as the one obtained by the Gilbert-Varshamov bound for standard offline binary error correcting codes, and in particular is optimal (up to constant factors). Interestingly, while it is known that the channel capacity in the online setting is $1-\Theta(\sqrt{\epsilon\log(1/\epsilon)})$--significantly lower than in the offline setting [20], the online requirement on the encoding function itself does not cost more in terms of the distance-rate trade-off. Rather, it is the additional overhead incurred by the mechanism required for synchronization that is responsible for the lower channel capacity in the online setting. We elaborate more on this in Section1.3.
The proof of Theorem4, which can be found in Section4, is based on a variant of the construction we use in Theorem2. There, the alphabet symbols are taken from the field of four elements, $\mathbb F_4$. The key idea in obtaining the savings in the alphabet size is to trace the $\mathbb F_4$ field elements down to $\mathbb F_2$ in most layers. Interestingly, we cannot afford to work over the field $\mathbb F_3$ as we crucially rely on the fact that the characteristic of the fields is2 as well as on the smaller field being a subfield of the larger one.
